% Figure: the barometric (gravitational Boltzmann) distribution. Relative
% number density of an isothermal ideal gas against altitude, for three gas
% species of different molar mass, showing that the heavier gas falls off
% faster with height. The density is n(z)/n(0) = exp(-m g z / k_B T), an
% instance of the Boltzmann factor with gravitational potential energy.
\begin{figure}[htbp]
\centering
\begin{tikzpicture}
\begin{axis}[
  width=12cm, height=7cm,
  xlabel={altitude $z$ (km)},
  ylabel={relative number density $n(z)/n(0)$},
  xmin=0, xmax=30, ymin=0, ymax=1.02,
  axis lines=left,
  legend pos=north east, legend cell align=left,
  legend style={font=\scriptsize, draw=enggray!40},
  tick label style={font=\scriptsize},
  label style={font=\small},
  smooth,
]
% scale height H = k_B T / (m g) in km, at T = 250 K.
% helium (M=4): H = R T /(M g) = 8.314*250/(0.004*9.81) = 52.9 km
\addplot[engblue, very thick, domain=0:30, samples=80]
  {exp(-x/52.9)};
\addlegendentry{helium ($M=4$, $H\approx53$ km)}
% nitrogen (M=28): H = 8.314*250/(0.028*9.81) = 7.56 km
\addplot[engred, very thick, domain=0:30, samples=80]
  {exp(-x/7.56)};
\addlegendentry{nitrogen ($M=28$, $H\approx7.6$ km)}
% carbon dioxide (M=44): H = 8.314*250/(0.044*9.81) = 4.81 km
\addplot[engamber, very thick, domain=0:30, samples=80]
  {exp(-x/4.81)};
\addlegendentry{carbon dioxide ($M=44$, $H\approx4.8$ km)}
\end{axis}
\end{tikzpicture}
\caption[The barometric distribution for three gas species]{The isothermal
barometric law $n(z)/n(0)=e^{-mgz/k_BT}$, the Boltzmann factor with
gravitational potential energy $mgz$ in place of an internal energy level,
drawn for three species at $T=250\,\si{\kelvin}$. Each curve falls off with a
scale height $H=k_BT/mg=RT/Mg$, the altitude over which the density drops by a
factor $e$; the heavier the molecule, the shorter the scale height and the
faster the density thins with altitude. Helium, light and weakly bound to the
planet, spreads through a scale height ten times that of carbon dioxide, which
hugs the ground. The same competition, thermal energy $k_BT$ pushing molecules
up against gravitational energy $mgz$ pulling them down, sets the equilibrium
profile: a balance struck not by force alone but by the entropy the gas gains
in spreading against the energy it costs to climb. The distribution is the
gravitational instance of the maximum-entropy Boltzmann form and underlies
sedimentation, centrifugal isotope separation, and the vertical structure of
every planetary atmosphere.}
\label{fig:vol04ch04:barometric-distribution}
\end{figure}
